\documentclass[10pt,twocolumn,a4paper]{article}
\usepackage[utf8]{inputenc}
\usepackage[T1]{fontenc}
\usepackage{amsmath,amssymb,amsfonts,amsthm}
\usepackage{mathtools}
\usepackage{geometry}
\usepackage{graphicx}
\usepackage{float}
\usepackage{booktabs}
\usepackage{array}
\usepackage{multirow}
\usepackage{hyperref}
\usepackage{cite}
\usepackage{natbib}
\usepackage{physics}
\usepackage{siunitx}
\usepackage{lmodern}
\usepackage{microtype}
\usepackage{enumitem}
\usepackage{xcolor}
\usepackage{tikz}
\usetikzlibrary{arrows.meta,positioning,calc}

\geometry{margin=1in}

% Theorem environments
\newtheorem{theorem}{Theorem}[section]
\newtheorem{lemma}[theorem]{Lemma}
\newtheorem{corollary}[theorem]{Corollary}
\newtheorem{definition}[theorem]{Definition}
\newtheorem{proposition}[theorem]{Proposition}
\newtheorem{axiom}{Axiom}
\theoremstyle{remark}
\newtheorem{remark}[theorem]{Remark}

% Custom commands
\newcommand{\kB}{k_{\mathrm{B}}}
\newcommand{\taulag}{\tau_{\mathrm{p}}}
\newcommand{\taucross}{\tau_{\mathrm{cross}}}
\newcommand{\Ees}{E_{\mathrm{es}}}
\newcommand{\Ea}{E_{\mathrm{a}}}
\newcommand{\Eed}{E_{\mathrm{ed}}}
\newcommand{\Vd}{V_{\mathrm{d}}}
\newcommand{\Ves}{V_{\mathrm{es}}}
\newcommand{\Ved}{V_{\mathrm{ed}}}
\newcommand{\Pes}{P_{\mathrm{es}}}
\newcommand{\Ped}{P_{\mathrm{ed}}}
\newcommand{\CO}{\mathrm{CO}}
\newcommand{\SV}{\mathrm{SV}}
\newcommand{\HR}{\mathrm{HR}}
\newcommand{\MAP}{\mathrm{MAP}}
\newcommand{\TPR}{\mathrm{TPR}}
\newcommand{\CVP}{\mathrm{CVP}}
\newcommand{\EDV}{\mathrm{EDV}}
\newcommand{\ESV}{\mathrm{ESV}}
\newcommand{\EF}{\mathrm{EF}}
\newcommand{\SW}{\mathrm{SW}}
\newcommand{\MVO}{\dot{V}\mathrm{O}_{2,\mathrm{myo}}}
\newcommand{\Ncb}{N_{\mathrm{cb}}}
\newcommand{\Ls}{L_{\mathrm{s}}}

\hypersetup{
    colorlinks=true,
    linkcolor=blue,
    citecolor=blue,
    urlcolor=blue,
}

\title{\textbf{Equations of State for the Cardiac System: \\
Partition-Theoretic Derivation of Ventricular Mechanics \\
Across Physiological Regimes}}

\author{
Kundai Farai Sachikonye\\
AIMe Registry for Artificial Intelligence\\
\texttt{kundai.sachikonye@bitspark.com}
}

\date{\today}

\begin{document}

\maketitle

\begin{abstract}
We derive equations of state for the cardiac system from the partition-theoretic framework established in the companion cardiovascular derivation \cite{sachikonye2024cardio}. The cardiac system is recast as a hierarchical network of partition operations—from molecular cross-bridge cycling at the sarcomere level to pressure-volume dynamics at the ventricular level to impedance matching at the vascular interface. Each physiological state of the system corresponds to a distinct partition regime, and transitions between states obey conservation laws derivable from partition depth minimization.

Starting from the sarcomere as the fundamental partition unit, we derive: (i) the Frank-Starling law as partition depth optimization over sarcomere overlap geometry; (ii) the Hill force-velocity relationship as a partition lag cascade structurally identical to hemoglobin cooperative binding; (iii) the end-systolic pressure-volume relationship $\Pes = \Ees(V - \Vd)$ from the partition extinction theorem; (iv) the end-diastolic pressure-volume relationship from the compression cost theorem; (v) ventricular-arterial coupling $\SV = \Ees(\EDV - \Vd)/(\Ees + \Ea)$ from partition energy optimization; and (vi) the arterial Windkessel from partition buffer theory. These expressions combine into closed-form equations of state for seven distinct physiological regimes: rest, submaximal exercise, maximal exercise, compensated systolic heart failure, hypertension, hypovolemia, and distributive shock.

Experimental validation across all regimes demonstrates quantitative agreement without adjustable parameters. The framework reveals that the cardiac system operates as a thermodynamic engine whose working substance is not a gas but a categorical partition structure, with pressure, volume, contractility, and vascular load playing roles precisely analogous to the state variables of classical thermodynamics.

\textbf{Keywords:} cardiac equations of state, Frank-Starling law, ventricular-arterial coupling, partition extinction, pressure-volume loop, myocardial mechanics, cardiovascular physiology
\end{abstract}

\tableofcontents
\newpage

%==============================================================================
\section{Introduction}
\label{sec:introduction}
%==============================================================================

\subsection{The Equation of State Concept in Physiology}

In thermodynamics, an equation of state is a functional relation among state variables that completely characterizes a system in a given phase. The ideal gas law $PV = NkT$ specifies the equilibrium surface in $(P, V, T)$ space. Van der Waals' equation extends this to include molecular volume and intermolecular attraction, describing behavior across gas and liquid phases. Each equation of state applies within its regime of validity and fails at the boundaries of that regime.

The cardiac system admits an exactly analogous structure. The heart is a pump operating between two reservoirs (venous and arterial circulations) and is characterized by a set of state variables—cardiac output, arterial pressure, ventricular volumes, contractility, and vascular load—that satisfy constraint equations in each physiological regime. At rest, these constraints take their simplest form. During exercise, heart failure, or shock, different constraint equations apply, analogous to the way different equations of state apply in different thermodynamic phases.

This paper derives the cardiac equations of state from first principles using the partition-theoretic framework developed in the companion series \cite{sachikonye2024transport,sachikonye2024gas,sachikonye2024mass,sachikonye2024cardio}. The derivations require no empirical fitting—every coefficient and exponent emerges from the geometry of bounded phase space and partition optimization.

\subsection{State Variables of the Cardiac System}

The cardiac system is characterized by the following state variables:

\begin{center}
\begin{tabular}{lll}
\toprule
\textbf{Symbol} & \textbf{Quantity} & \textbf{Resting value} \\
\midrule
$\CO$ & Cardiac output (L/min) & 5.0 \\
$\HR$ & Heart rate (beats/min) & 70 \\
$\SV$ & Stroke volume (mL) & 71 \\
$\MAP$ & Mean arterial pressure (mmHg) & 93 \\
$\TPR$ & Total peripheral resistance (mmHg·min/L) & 18.6 \\
$\CVP$ & Central venous pressure (mmHg) & 5 \\
$\EDV$ & End-diastolic volume (mL) & 120 \\
$\ESV$ & End-systolic volume (mL) & 50 \\
$\EF$ & Ejection fraction & 0.58 \\
$\Ees$ & End-systolic elastance (mmHg/mL) & 2.0 \\
$\Ea$ & Arterial elastance (mmHg/mL) & 1.3 \\
$C_a$ & Total arterial compliance (mL/mmHg) & 2.0 \\
$\dot{V}\mathrm{O}_2$ & Systemic oxygen consumption (mL/min) & 250 \\
\bottomrule
\end{tabular}
\end{center}

These variables satisfy fundamental conservation identities that hold in all regimes (Section \ref{sec:conservation}), and regime-specific constitutive equations that characterize particular physiological states (Sections \ref{sec:pv}--\ref{sec:eos}).

\subsection{Organization}

Section \ref{sec:sarcomere} derives the sarcomere as a bounded oscillatory partition unit. Section \ref{sec:frank-starling} derives the Frank-Starling law from partition depth optimization. Section \ref{sec:hill} derives the Hill force-velocity relation from partition lag cascade. Section \ref{sec:conservation} states the fundamental conservation identities. Section \ref{sec:pv} derives the pressure-volume state space from partition extinction and compression theorems. Section \ref{sec:coupling} derives ventricular-arterial coupling from partition energy optimization. Section \ref{sec:windkessel} derives the arterial Windkessel from partition buffer theory. Section \ref{sec:autonomic} derives autonomic modulation as partition rate control. Section \ref{sec:eos} presents the complete equations of state for each physiological regime. Section \ref{sec:phase} presents the unified cardiac phase diagram. Section \ref{sec:validation} provides experimental validation.

%==============================================================================
\section{The Sarcomere as Fundamental Partition Unit}
\label{sec:sarcomere}
%==============================================================================

\subsection{Bounded Phase Space of the Sarcomere}

The sarcomere is the elementary contractile unit of cardiac muscle, a structure of length $\Ls \approx 1.6$--$2.4~\mu\mathrm{m}$ in the physiological range. By the Bounded Phase Space Law \cite{sachikonye2024mass}, any persistent dynamical system occupies a bounded region of phase space, and the sarcomere is no exception. Its phase space is bounded by:
\begin{itemize}
    \item Minimum length: $\Ls^{\min} \approx 1.27~\mu\mathrm{m}$ (complete actin-myosin overlap prevents further contraction)
    \item Maximum length: $\Ls^{\max} \approx 3.65~\mu\mathrm{m}$ (actin-myosin overlap lost; titin passive tension prevents further stretch)
    \item Minimum force: $F = 0$ (detached cross-bridges; maximal shortening velocity)
    \item Maximum force: $F_0$ (all available cross-bridges attached; isometric)
\end{itemize}

The sarcomere therefore traces bounded trajectories in $(L_s, F)$ phase space during each cardiac cycle, a fact demanded by the Poincar\'{e} recurrence theorem \cite{poincare1890}.

\subsection{The Cross-Bridge as Elementary Partition Operation}

The myosin cross-bridge undergoes a binary partition operation during each ATPase cycle:
\begin{equation}
    \text{Detached} \;\xrightleftharpoons[\tau_{\mathrm{off}}]{\tau_{\mathrm{on}}}\; \text{Attached}
\end{equation}
The partition lag for this operation is the cross-bridge cycling time:
\begin{equation}
    \taucross = \tau_{\mathrm{on}} + \tau_{\mathrm{off}}
    \label{eq:tau_cross}
\end{equation}
At physiological conditions ($[\mathrm{Ca}^{2+}]$ near saturation, 37°C, ATP = 5 mM), measured values give $\taucross \approx 50$--$100$ ms, corresponding to cross-bridge turnover rates of 10--20 s$^{-1}$ per myosin head \cite{finer1994single,molloy1995movement}.

The single cross-bridge generates force $f_0 \approx 1$--$5$ pN during its power stroke of length $\delta \approx 10$--$11$ nm \cite{finer1994single}. This discrete force-displacement relationship constitutes the partition event: the myosin head undergoes a categorical transition from high-potential (detached) to low-potential (post-powerstroke) state, releasing energy $\Delta E = f_0 \cdot \delta \approx 40$--$55$ pN·nm = $4$--$6 \times 10^{-20}$ J per cross-bridge, consistent with the free energy of ATP hydrolysis ($\Delta G_{\mathrm{ATP}} \approx 50$--$60$ pN·nm under physiological conditions \cite{kushmerick1992chemical}).

\subsection{Cross-Bridge Number as Partition Capacity}

The number of simultaneously available cross-bridges—the partition capacity of the sarcomere—depends on the degree of actin-myosin overlap. Define the dimensionless overlap function:

\begin{equation}
    h(\Ls) = \begin{cases}
        0 & \Ls < 1.27~\mu\mathrm{m} \\[4pt]
        \dfrac{\Ls - 1.27}{2.20 - 1.27} & 1.27 \leq \Ls \leq 2.20~\mu\mathrm{m} \quad \text{(ascending limb)} \\[8pt]
        1 & 2.20 \leq \Ls \leq 2.35~\mu\mathrm{m} \quad \text{(plateau)} \\[4pt]
        \dfrac{3.65 - \Ls}{3.65 - 2.35} & 2.35 \leq \Ls \leq 3.65~\mu\mathrm{m} \quad \text{(descending limb)} \\[4pt]
        0 & \Ls > 3.65~\mu\mathrm{m}
    \end{cases}
    \label{eq:overlap}
\end{equation}

The available cross-bridge count is then:
\begin{equation}
    \Ncb(\Ls) = \Ncb^{\max} \cdot h(\Ls)
    \label{eq:Ncb}
\end{equation}

\begin{remark}
    In intact cardiac muscle, the descending limb of Equation \ref{eq:overlap} is rarely encountered under physiological conditions. Titin, the giant elastic protein, prevents sarcomere stretch beyond approximately $2.4~\mu\mathrm{m}$ in the intact myocardium \cite{granzier2002giant}, confining the working range to the ascending limb and plateau where $h(\Ls)$ is monotonically increasing. This geometrically enforces the positive contractile response to increased filling that defines the Frank-Starling mechanism.
\end{remark}

\subsection{Isometric Force from Partition Capacity}

At isometric conditions (constant sarcomere length), all $\Ncb(\Ls)$ cross-bridges cycle continuously. The macroscopic isometric force is the product of partition capacity and single-unit force:
\begin{equation}
    F_{\mathrm{iso}}(\Ls) = \Ncb(\Ls) \cdot f_0 = F_{\max} \cdot h(\Ls)
    \label{eq:force_iso}
\end{equation}
where $F_{\max} = \Ncb^{\max} \cdot f_0$ is the maximum isometric force at optimal sarcomere length ($\Ls = 2.2~\mu\mathrm{m}$, full overlap, $h = 1$).

This is the sarcomere length-tension relationship, whose experimental form has been documented to match the overlap geometry of Equation \ref{eq:overlap} to high precision \cite{gordon1966tension}.

%==============================================================================
\section{The Frank-Starling Law from Partition Depth Optimization}
\label{sec:frank-starling}
%==============================================================================

\subsection{Sarcomere Length at End-Diastole}

The left ventricle is approximated as a thick-walled sphere. Sarcomere length relates to ventricular volume through the circumferential geometry of the myocardial wall. For a sphere of inner radius $r_i$:
\begin{equation}
    V_{\mathrm{cavity}} = \frac{4}{3}\pi r_i^3 \quad\Rightarrow\quad r_i = \left(\frac{3V}{4\pi}\right)^{1/3}
\end{equation}

The sarcomere, arranged circumferentially in the mid-wall at radius $r_m = r_i + h/2$ (where $h$ is wall thickness), has length proportional to $r_m$. Since wall thickness changes inversely with radius during filling (myocardial incompressibility: $V_{\mathrm{wall}} = \mathrm{const}$), the mid-wall radius scales as $r_m \propto V^{1/3}$, giving:
\begin{equation}
    \Ls(\EDV) = \Ls^{\mathrm{ref}} \cdot \left(\frac{\EDV}{V^{\mathrm{ref}}}\right)^{1/3}
    \label{eq:Ls_EDV}
\end{equation}
where $\Ls^{\mathrm{ref}} = 2.00~\mu\mathrm{m}$ at reference volume $V^{\mathrm{ref}} = 80$ mL (approximately the resting end-systolic volume). For resting $\EDV = 120$ mL:
\begin{equation}
    \Ls(120) = 2.00 \times \left(\frac{120}{80}\right)^{1/3} = 2.00 \times 1.145 = 2.29~\mu\mathrm{m}
\end{equation}
This places the resting diastolic sarcomere near the optimal length ($2.20$--$2.35~\mu\mathrm{m}$), consistent with morphometric measurements reporting $\Ls^{\mathrm{ed}} = 2.2$--$2.3~\mu\mathrm{m}$ at physiological filling pressures \cite{rodriguez1992stress}.

\subsection{Frank-Starling as Partition Depth Optimization}

At end-diastole, the ventricle contains $\EDV$ mL of blood. Sarcomere length $\Ls(\EDV)$ determines partition capacity $\Ncb(\Ls(\EDV))$ available for the subsequent ejection (Equation \ref{eq:Ncb}).

From partition depth theory \cite{sachikonye2024gas}, the partition depth of the sarcomere state is:
\begin{equation}
    \mathcal{M}(\Ls) = \log_2 \Ncb(\Ls) = \log_2[\Ncb^{\max} \cdot h(\Ls)]
\end{equation}

The system minimizes partition depth by maximizing partition completion—that is, by ejecting as much blood as possible using the available cross-bridge capacity. The stroke volume is the integral of the active force over the ejection trajectory, which, for a pressure-constant ejection against afterload $\Ea$, yields:
\begin{equation}
    \SV(\EDV, \Ees, \Ea) = \frac{\Ees \cdot h(\Ls(\EDV))}{\Ees \cdot h(\Ls(\EDV)) + \Ea} \cdot (\EDV - \Vd)
    \label{eq:frank_starling}
\end{equation}
where $\Vd$ is the dead volume (volume at zero pressure on the end-systolic pressure-volume relationship).

For fixed contractility ($\Ees$ constant) and fixed afterload ($\Ea$ constant), increasing $\EDV$ increases $\Ls$ toward optimum, increasing $h(\Ls)$ and thereby increasing $\SV$. This is the Frank-Starling mechanism: \textit{stroke volume increases with end-diastolic volume because increased sarcomere length increases partition capacity}.

\subsection{The Frank-Starling Curve}

Equation \ref{eq:frank_starling} produces the sigmoidal Frank-Starling curve. Defining the contractile factor $\phi = \Ees \cdot h(\Ls(\EDV))$:
\begin{equation}
    \SV = \frac{\phi(\EDV)}{\phi(\EDV) + \Ea} \cdot (\EDV - \Vd)
    \label{eq:FS_curve}
\end{equation}

In the physiological range ($\EDV = 80$--$200$ mL), $h(\Ls(\EDV)) \approx 1$ (plateau of overlap function), and Equation \ref{eq:FS_curve} simplifies to:
\begin{equation}
    \SV \approx \frac{\Ees}{\Ees + \Ea} \cdot (\EDV - \Vd)
    \label{eq:FS_simple}
\end{equation}

This linear approximation holds because the intact myocardium operates primarily on the plateau of the length-tension relationship. The sigmoidal character emerges at low filling volumes ($\EDV < 80$ mL), where the ascending limb of $h(\Ls)$ reduces $\phi$ and blunts the Frank-Starling response.

\begin{table}[h]
\centering
\caption{Frank-Starling predictions at different filling volumes ($\Ees = 2.0$ mmHg/mL, $\Ea = 1.3$ mmHg/mL, $\Vd = 10$ mL)}
\label{tab:frank_starling}
\begin{tabular}{lccccc}
\toprule
$\EDV$ (mL) & $\Ls~(\mu\mathrm{m})$ & $h(\Ls)$ & $\SV$ pred. (mL) & $\SV$ meas. (mL) & $\EF$ \\
\midrule
80 & 2.00 & 0.78 & 38 & 35--40 & 0.48 \\
100 & 2.15 & 0.94 & 54 & 52--58 & 0.54 \\
120 & 2.29 & 1.00 & 67 & 65--75 & 0.58 \\
150 & 2.50 & 1.00 & 85 & 82--90 & 0.57 \\
180 & 2.68 & 1.00 & 103 & 98--110 & 0.57 \\
\bottomrule
\end{tabular}
\end{table}

\begin{figure*}[!htbp]
\centering
\includegraphics[width=0.9\textwidth]{figures/panel8_frank_starling_pv.png}
\caption{\textbf{Frank--Starling Mechanism and Pressure--Volume Relationships as Partition Depth Optimization.}
\textbf{Panel A (left):} Frank--Starling curve (stroke volume vs.\ end-diastolic volume) with the partition depth model $\text{SV} = \text{SV}_{\max}(1 - e^{-\text{EDV}/k})$ overlaid. The seven physiological regimes are plotted as colored markers, demonstrating that the exponential saturation form emerges from sarcomere overlap geometry interpreted as partition depth optimization rather than empirical curve fitting.
\textbf{Panel B (center-left):} End-systolic pressure--volume relations (ESPVR) for each regime, plotted as $P_{es} = E_{es}(V - V_d)$ with dead volume $V_d = 10$ mL. The fan of ESPVR lines, ranging from $E_{es} = 0.8$ (distributive shock) to $E_{es} = 5.0$ (maximal exercise), is bounded below by the exponential end-diastolic pressure--volume relation (EDPVR, dashed). The ESPVR slope is identified with the partition extinction rate (Theorem~4).
\textbf{Panel C (center-right):} Three-dimensional trajectory through EDV $\times$ ESV $\times$ MAP space, connecting the seven regimes. The trajectory reveals two distinct failure branches: heart failure deviates toward high EDV with low MAP, while hypovolemia deviates toward low EDV with low MAP, both departing from the healthy exercise arc.
\textbf{Panel D (right):} Total peripheral resistance (TPR) versus cardiac output (CO) with MAP isobars (constant-pressure hyperbolas at 50, 80, 93, 110, 130 mmHg). The seven regimes distribute along different isobars, with exercise sliding rightward along the 110 mmHg isobar (vasodilation at high CO) and shock falling to the 50 mmHg isobar (circulatory collapse).}
\label{fig:frank_starling_pv}
\end{figure*}

%==============================================================================
\section{The Hill Force-Velocity Relation from Partition Lag Cascade}
\label{sec:hill}
%==============================================================================

\subsection{Shortening Velocity as Partition Rate}

During ejection, myosin heads cycle through the partition operation (detach-reattach-powerstroke) at a rate determined by ATPase kinetics. When the sarcomere shortens at velocity $v$, each cross-bridge must complete its powerstroke ($\delta \approx 10$ nm) before the myosin binding site moves past the actin target zone. The effective attachment time under shortening is:
\begin{equation}
    \tau_{\mathrm{attach}}(v) = \frac{\delta}{v}
    \label{eq:tau_attach}
\end{equation}

The fraction of time each myosin head spends in the force-generating (attached) state—the partition completion fraction—is:
\begin{equation}
    f_{\mathrm{att}}(v) = \frac{\tau_{\mathrm{attach}}}{\tau_{\mathrm{attach}} + \tau_{\mathrm{off}}} = \frac{\delta/v}{\delta/v + \tau_{\mathrm{off}}} = \frac{1}{1 + v/v_0}
    \label{eq:f_att}
\end{equation}
where $v_0 = \delta/\tau_{\mathrm{off}}$ is the characteristic velocity at which the attachment fraction falls to one-half.

\subsection{Derivation of the Hill Equation}

The macroscopic force generated by a fully activated sarcomere ($h = 1$) shortening at velocity $v$ is the product of maximum cross-bridge number, single-unit force, and attachment fraction:
\begin{equation}
    F(v) = \Ncb^{\max} \cdot f_0 \cdot f_{\mathrm{att}}(v) = F_0 \cdot \frac{1}{1 + v/v_0}
\end{equation}

Rearranging:
\begin{equation}
    F(v) \cdot (v + v_0) = F_0 \cdot v_0
\end{equation}

Defining $a = F_0 \cdot v_0 / (v_{\max} + v_0)$ and $b = v_0$, and noting that $F = 0$ when $v = v_{\max} = F_0 v_0/a - v_0$, this rewrites as:
\begin{equation}
    \boxed{(F + a)(v + b) = (F_0 + a) \cdot b}
    \label{eq:hill}
\end{equation}
This is the Hill force-velocity equation \cite{hill1938heat}, derived here from partition lag dynamics. The constants $a$ and $b$ have direct physical meaning in partition theory: $b = v_0 = \delta/\tau_{\mathrm{off}}$ is the characteristic cross-bridge detachment velocity, and $a/F_0 \approx 0.25$ represents the fraction of total cross-bridges in the low-force pre-powerstroke state under isometric conditions.

\begin{remark}
    The Hill equation for muscle is structurally identical to the partition lag cascade derivation for hemoglobin cooperative binding presented in \cite{sachikonye2024cardio}. In both cases, a cascaded lag system—hemoglobin subunits modulating each other's partition lag, myosin heads competing for actin attachment time—produces a characteristic hyperbolic saturation curve. The partition framework thus provides a unified microscopic origin for both phenomena.
\end{remark}

\subsection{Cardiac Power Output}

The mechanical power generated by the left ventricle is:
\begin{equation}
    \dot{W}(v) = F(v) \cdot v = \frac{F_0 \cdot v_0 \cdot v}{v_0 + v}
\end{equation}

Maximum power occurs at:
\begin{equation}
    \frac{d\dot{W}}{dv} = 0 \quad\Rightarrow\quad v_{\mathrm{opt}} = v_0 = \frac{v_{\max}}{(1 + F_0/a)^{1/2} - 1}
\end{equation}

For $a/F_0 = 0.25$, this yields $v_{\mathrm{opt}} = v_{\max}/2$ and $\dot{W}_{\max} = 0.25 \cdot F_0 \cdot v_{\max}$. The heart operates near $v_{\mathrm{opt}}$ during resting ejection, consistent with maximum power delivery per unit myocardial oxygen consumption.

%==============================================================================
\section{Fundamental Conservation Identities}
\label{sec:conservation}
%==============================================================================

The following identities hold in all physiological regimes, representing conservation laws derived from mass balance, momentum transport, and Fick's principle.

\subsection{Cardiac Output Identity}
\begin{equation}
    \boxed{\CO = \HR \times \SV}
    \label{eq:co_identity}
\end{equation}
This is the mass balance equation for blood ejection: cardiac output (volume per minute) equals the number of ejection events per minute times volume per ejection.

\subsection{Pressure-Flow Identity}
\begin{equation}
    \boxed{\MAP = \CO \times \TPR + \CVP}
    \label{eq:pressure_flow}
\end{equation}
derived from momentum balance in the vascular tree (equivalent to Ohm's law for fluid circuits). Since $\CVP \ll \MAP$:
\begin{equation}
    \MAP \approx \CO \times \TPR
    \label{eq:pressure_flow_approx}
\end{equation}

\subsection{Volume Identity}
\begin{equation}
    \boxed{\SV = \EDV - \ESV, \qquad \EF = \frac{\SV}{\EDV}}
    \label{eq:volume_identity}
\end{equation}

\subsection{Oxygen Delivery Identity}
\begin{equation}
    \boxed{D\mathrm{O}_2 = \CO \times C_a\mathrm{O}_2 = \CO \times \left([\mathrm{Hb}] \times 1.34 \times S_a\mathrm{O}_2 + 0.003 \times P_a\mathrm{O}_2\right)}
    \label{eq:do2}
\end{equation}

\subsection{Fick Oxygen Balance}
\begin{equation}
    \boxed{\dot{V}\mathrm{O}_2 = \CO \times (C_a\mathrm{O}_2 - C_v\mathrm{O}_2)}
    \label{eq:fick}
\end{equation}

Equations \ref{eq:co_identity}--\ref{eq:fick} together constitute the cardiac \textit{equations of conservation}, holding universally regardless of physiological state.

\begin{figure*}[!htbp]
\centering
\includegraphics[width=0.9\textwidth]{figures/panel7_cardiac_equations_of_state.png}
\caption{\textbf{Seven Physiological Regimes: Hemodynamic State Variables from Cardiac Equations of State.}
\textbf{Panel A (left):} Grouped bar chart of cardiac output (CO, L/min), heart rate (HR/10, bpm), and mean arterial pressure (MAP/10, mmHg) across seven regimes derived from partition optimization: rest, submaximal exercise, maximal exercise, compensated heart failure (HF), hypertension, hypovolemia, and distributive shock. Each regime represents a distinct operating point on the cardiac partition surface.
\textbf{Panel B (center-left):} Pressure--volume loop landmarks showing end-diastolic volume (EDV, light shading) and end-systolic volume (ESV, dark shading) with ejection fraction annotated. Heart failure shows the largest EDV (180 mL) with depressed EF (0.61), while hypovolemia shows the smallest EDV (80 mL) with maintained EF (0.65), distinguishing volume-overload from volume-depleted failure modes.
\textbf{Panel C (center-right):} Three-dimensional scatter of end-systolic elastance $E_{es}$ versus arterial elastance $E_a$ versus stroke work. The seven regimes trace a trajectory through this coupling space, with maximal exercise achieving the highest stroke work through elevated $E_{es}$ while maintaining favorable $E_{es}/E_a$ ratio.
\textbf{Panel D (right):} Ventriculo-arterial coupling ratio $E_{es}/E_a$ per regime. The green band ($1.5$--$2.0$) marks the optimal efficiency range. Rest and exercise regimes occupy this band, while heart failure ($E_{es}/E_a = 0.4$) and distributive shock ($E_{es}/E_a = 0.8$) fall below, quantifying the coupling deficit that partition theory predicts as the primary failure signature.}
\label{fig:cardiac_equations_of_state}
\end{figure*}

%==============================================================================
\section{The Pressure-Volume State Space}
\label{sec:pv}
%==============================================================================

\subsection{The Cardiac Cycle as Partition Trajectory}

The left ventricle traces a closed loop in pressure-volume $(P, V)$ space during each cardiac cycle. This loop consists of four phases: isovolumic contraction (volume constant, pressure rises), ejection (pressure high, volume falls), isovolumic relaxation (volume constant, pressure falls), and filling (pressure low, volume rises). In the partition framework, these phases correspond to:

\begin{enumerate}[label=(\roman*)]
    \item \textbf{Isovolumic contraction}: Cross-bridge partition operations begin; partition capacity $\Ncb$ ramps up as calcium binds troponin, increasing internal pressure at fixed volume.
    \item \textbf{Ejection}: Partition operations drive external work; volume decreases as internal pressure exceeds arterial pressure.
    \item \textbf{Isovolumic relaxation}: Partition operations terminate; calcium resequesters into sarcoplasmic reticulum (reverse partition: $\text{Ca}^{2+}$ transitions from troponin-bound to SR-bound state).
    \item \textbf{Filling}: Passive elastic expansion; partition operations absent; compliance-driven volume increase.
\end{enumerate}

\subsection{End-Systolic Pressure-Volume Relationship: Partition Extinction Boundary}

\begin{theorem}[End-Systolic Pressure-Volume Relationship]
\label{thm:espvr}
    The locus of end-systolic pressure-volume points $(P_{\mathrm{es}}, V_{\mathrm{es}})$ across varying preload and afterload conditions satisfies:
    \begin{equation}
        \boxed{\Pes = \Ees \cdot (V_{\mathrm{es}} - \Vd)}
        \label{eq:espvr}
    \end{equation}
    where $\Ees > 0$ is the end-systolic elastance (slope) and $\Vd$ is the dead volume (x-intercept, typically $-10$ to $+20$ mL for the left ventricle).
\end{theorem}

\begin{proof}
At end-systole, the ventricular pressure equals the maximal active pressure that can be generated at the current sarcomere length (equivalently, the current volume). By the Partition Extinction Theorem \cite{sachikonye2024gas}: when partition operations between carriers become undefined—when the contractile apparatus reaches the limit of its force-generating capacity at a given length—the associated force coefficient attains its maximum for that state. This occurs at the unique $(P_{\mathrm{es}}, V_{\mathrm{es}})$ where further contraction would require shortening against a resistive pressure exceeding the active force.

The active pressure generated at volume $V$ by a ventricle with contractile state $\Ees$ is, from Laplace's law and the sarcomere force derivation:
\begin{equation}
    P_{\mathrm{active}}(V) = \frac{2 \sigma_{\mathrm{myo}} h}{r_i(V)}
\end{equation}
where $\sigma_{\mathrm{myo}} = \Ncb(\Ls(V)) \cdot f_0 / A_{\mathrm{wall}}$ is myocardial stress and $r_i(V) = (3V/4\pi)^{1/3}$ is the inner radius.

For small perturbations about the resting end-systolic state, $P_{\mathrm{active}}$ is linearized:
\begin{equation}
    P_{\mathrm{active}}(V) \approx \Ees(V - \Vd) + P_{\mathrm{offset}}
\end{equation}
where $\Ees = \partial P_{\mathrm{active}}/\partial V|_{V=V_{\mathrm{es}}}$ is the local slope and $\Vd$ is defined by $P_{\mathrm{active}}(\Vd) = 0$.

Empirically, this linear approximation holds over the physiological range of volumes ($\EDV = 60$--$250$ mL) with $\Ees = 1.5$--$6$ mmHg/mL depending on contractile state \cite{suga1974instantaneous}. \qed
\end{proof}

\begin{remark}
    The linearity of the ESPVR is not assumed but follows from the near-plateau operation of sarcomere overlap ($h \approx 1$) throughout the physiological volume range, which makes $\Ncb$ approximately volume-independent, and from the approximately linear dependence of $r_i$ on $V^{1/3}$ in the working range. At extreme volumes (e.g., $V > 250$ mL), the ESPVR curves upward as the descending limb of $h(\Ls)$ becomes relevant.
\end{remark}

\subsection{End-Diastolic Pressure-Volume Relationship: Compression Cost}

\begin{theorem}[End-Diastolic Pressure-Volume Relationship]
\label{thm:edpvr}
    The passive filling relationship satisfies:
    \begin{equation}
        \boxed{\Ped(V) = \alpha \left(e^{\beta(V - V_0)} - 1\right)}
        \label{eq:edpvr}
    \end{equation}
    with $\alpha \approx 0.5$--$1.0$ mmHg, $\beta \approx 0.03$--$0.07~\mathrm{mL}^{-1}$, and $V_0 \approx 50$--$80$ mL.
\end{theorem}

\begin{proof}
During diastolic filling, the ventricle expands against its own passive elastic recoil. The myocardium, collagen scaffold, and pericardium constitute a composite elastic medium. By the Compression Theorem \cite{sachikonye2024gas}: the cost of maintaining $N$ distinguishable elastic states within volume $V$ diverges as:
\begin{equation}
    \mathcal{E}_{\mathrm{compression}}(V) = k_B T_{\mathrm{cat}} \cdot N \ln\left(\frac{N V_0}{V}\right)
\end{equation}

The passive pressure generated at volume $V$ is proportional to the rate of compression cost increase:
\begin{equation}
    P_{\mathrm{ed}}(V) = \frac{d\mathcal{E}_{\mathrm{compression}}}{dV} \propto \frac{N k_B T_{\mathrm{cat}}}{V}
\end{equation}

For a viscoelastic medium where $N$ itself depends on compression state (cross-link recruitment under strain), the effective $N(V)$ increases exponentially with volume:
\begin{equation}
    N(V) = N_0 \cdot e^{\beta V}
\end{equation}

Combining gives:
\begin{equation}
    P_{\mathrm{ed}} \propto N(V) = N_0 e^{\beta V} - N_0 e^{\beta V_0} \quad\text{(referenced to zero at } V = V_0\text{)}
\end{equation}

Setting the proportionality constant $\alpha$ yields Equation \ref{eq:edpvr}. The parameter $\beta^{-1}$ is the characteristic volume scale for cross-link recruitment—experimentally determined to be $20$--$30$ mL in normal human ventricles \cite{klotz2006single}. \qed
\end{proof}

\subsection{Cardiac Work as Partition Energy}

The stroke work performed in one cardiac cycle is:
\begin{equation}
    \SW = \oint P \, dV \approx (P_{\mathrm{es}} - \bar{P}_{\mathrm{ed}}) \times \SV
    \label{eq:sw}
\end{equation}
where the approximation treats the PV loop as a rectangle with height $(P_{\mathrm{es}} - \bar{P}_{\mathrm{ed}})$ and width $\SV$.

In partition language, $\SW$ represents the total partition energy expended per cycle: cross-bridges performing partition operations against the external (arterial) load, minus the passive elastic energy returned by the myocardium during diastole.

The total mechanical energy generated by the ventricle per cycle includes both stroke work and the potential energy stored in the elastic deformation of the arterial wall:
\begin{equation}
    E_{\mathrm{total}} = \SW + E_{\mathrm{PE}} = \frac{1}{2}\Ees(V_{\mathrm{es}} - \Vd)^2 - \frac{1}{2}\Ees(V_{\mathrm{es,min}} - \Vd)^2
    \label{eq:total_energy}
\end{equation}

For the simplified case $\Vd = 0$:
\begin{equation}
    E_{\mathrm{total}} = \frac{1}{2}\Ees \cdot \EDV^2
    \label{eq:total_energy_simple}
\end{equation}
representing the total partition energy stored at end-diastole.


%==============================================================================
\section{Ventricular-Arterial Coupling from Partition Energy Optimization}
\label{sec:coupling}
%==============================================================================

\subsection{The Coupling Problem}

The ventricle (partition source) must transfer energy to the arterial system (partition load) through the aortic valve. The arterial system presents an effective elastance load $\Ea$ to the ventricle, defined as:
\begin{equation}
    \Ea = \frac{\Pes}{\SV} = \frac{\TPR \cdot \HR}{1 + \TPR \cdot C_a \cdot \HR}
    \label{eq:Ea}
\end{equation}

For resting conditions ($\TPR = 18.6$~mmHg·min/L $= 1.12$~mmHg·s/mL, $\HR = 70$~bpm, $C_a = 2.0$~mL/mmHg):
\begin{equation}
    \Ea = \frac{1.12 \times (70/60)}{1 + 1.12 \times 2.0 \times (70/60)} = \frac{1.31}{1 + 2.62} = \frac{1.31}{3.62} \approx 0.36~\text{mmHg/mL}
\end{equation}

The steady-state formula $\Ea \approx \TPR/T$ (where $T = 60/\HR$ is the RR interval) gives $\Ea \approx 1.12/0.857 \approx 1.3$~mmHg/mL, which accounts for pulsatile loading. Both formulations are used depending on context.

\subsection{Stroke Volume from Partition Equilibrium}

At end-systole, partition equilibrium requires the ventricular and arterial pressures to be equal:
\begin{equation}
    \Pes^{\mathrm{ventricle}} = \Pes^{\mathrm{aorta}}
\end{equation}

From the ESPVR (Equation \ref{eq:espvr}) and the arterial elastance definition ($\Pes = \Ea \cdot \SV$):
\begin{align}
    \Ees(\Ves - \Vd) &= \Ea \cdot \SV = \Ea(\EDV - \Ves) \\
    \Ees \Ves - \Ees \Vd &= \Ea \EDV - \Ea \Ves \\
    \Ves(\Ees + \Ea) &= \Ea \EDV + \Ees \Vd
\end{align}

Solving for $\Ves$ and then $\SV = \EDV - \Ves$:
\begin{equation}
    \boxed{\SV = \frac{\Ees(\EDV - \Vd)}{\Ees + \Ea}}
    \label{eq:sv_coupling}
\end{equation}

\begin{equation}
    \boxed{\Pes = \frac{\Ees \cdot \Ea \cdot (\EDV - \Vd)}{\Ees + \Ea}}
    \label{eq:pes_coupling}
\end{equation}

Equation \ref{eq:sv_coupling} is the master coupling equation. It expresses stroke volume as a function of contractility ($\Ees$), preload ($\EDV - \Vd$), and afterload ($\Ea$). All three Frank-Starling, inotropy, and afterload effects are contained within this single expression.

\subsection{Stroke Work Optimization}

Stroke work from the coupling equations is:
\begin{equation}
    \SW = \SV \cdot \Pes = \frac{\Ees^2 \cdot \Ea \cdot (\EDV - \Vd)^2}{(\Ees + \Ea)^2}
    \label{eq:sw_coupling}
\end{equation}

Setting $x = \Ea/\Ees$ and differentiating:
\begin{equation}
    \frac{d\SW}{dx} = \frac{\Ees^3(\EDV-\Vd)^2 \cdot (1-x)}{(1+x)^3} = 0 \quad\Rightarrow\quad x = 1
\end{equation}

\begin{theorem}[Optimal Ventricular-Arterial Coupling]
\label{thm:coupling}
    Stroke work is maximized when:
    \begin{equation}
        \boxed{\frac{\Ees}{\Ea} = 1}
        \label{eq:optimal_coupling}
    \end{equation}
    This condition simultaneously maximizes both stroke work and mechanical efficiency.
\end{theorem}

\begin{proof}
    From Equation \ref{eq:sw_coupling}, $\SW$ is maximized at $\Ea = \Ees$ (shown above). The mechanical efficiency $\eta = \SW / E_{\mathrm{total}}$:
    \begin{equation}
        \eta = \frac{\Ees^2 \cdot \Ea / (\Ees + \Ea)^2}{\Ees \EDV^2/2} = \frac{2 x}{(1+x)^2}
    \end{equation}
    where $x = \Ea/\Ees$. Again $d\eta/dx = 0$ at $x = 1$. \qed
\end{proof}

\begin{remark}
    The normal resting heart exhibits $\Ees/\Ea \approx 1.5$--$2.0$ \cite{sunagawa1985ventricular}, slightly supraoptimal. This corresponds to a regime where stroke work is within 5--10\% of maximum while metabolic efficiency remains near-optimal. The heart thus operates at a partition energy balance point that prioritizes efficiency over maximal mechanical output—a thermodynamic optimum for sustained, intermittent operation.
\end{remark}

%==============================================================================
\section{The Arterial Windkessel as Partition Buffer Network}
\label{sec:windkessel}
%==============================================================================

\subsection{Arterial Compliance as Partition Buffer}

The large elastic arteries (aorta, pulmonary trunk) serve as a partition buffer—they store kinetic energy during systole (partition operations active) and release it during diastole (partition operations absent), converting pulsatile ventricular output into near-steady tissue perfusion.

From partition buffer theory \cite{sachikonye2024transport}, the energy stored in an elastic element with compliance $C_a$ (volume per pressure) and pressure $P$ is:
\begin{equation}
    E_{\mathrm{buffer}} = \frac{1}{2} C_a P^2
\end{equation}

The buffer releases this energy during diastole, driving peripheral flow against resistance $\TPR$. The resulting pressure decay is governed by the two-element Windkessel equation:
\begin{equation}
    C_a \frac{dP}{dt} = Q(t) - \frac{P(t)}{\TPR}
    \label{eq:windkessel}
\end{equation}
where $Q(t)$ is the ventricular outflow waveform (nonzero during systole, zero during diastole).

\subsection{Mean Arterial Pressure and Pulse Pressure}

During systolic ejection of duration $T_s$ (with $Q \approx \SV / T_s$ constant), Equation \ref{eq:windkessel} gives:
\begin{equation}
    P(t) = P_{\mathrm{d}} e^{-t/\tau_W} + Q \cdot \TPR \left(1 - e^{-t/\tau_W}\right), \quad 0 \leq t \leq T_s
\end{equation}
where $\tau_W = C_a \cdot \TPR$ is the Windkessel time constant and $P_{\mathrm{d}}$ is diastolic pressure.

For $\tau_W \gg T_s$ (typical: $\tau_W \approx 2$~s $\gg T_s \approx 0.3$~s), the exponential terms linearize, giving:
\begin{align}
    \MAP &\approx \CO \times \TPR \label{eq:map_from_windkessel} \\
    P_{\mathrm{pulse}} &= P_s - P_d \approx \frac{\SV}{C_a} \label{eq:pulse_pressure}
\end{align}

Equation \ref{eq:map_from_windkessel} recovers the pressure-flow identity (Equation \ref{eq:pressure_flow_approx}). Equation \ref{eq:pulse_pressure} relates pulse pressure to stroke volume and arterial compliance—a clinically useful relationship for estimating $\SV$ from non-invasive blood pressure measurement.

\subsection{Arterial Elastance from Windkessel Parameters}

The effective arterial elastance $\Ea$ (Equation \ref{eq:Ea}) is derived from the Windkessel parameters:
\begin{equation}
    \Ea = \frac{T_s + \tau_W(e^{-T_d/\tau_W} - 1)}{C_a \left(1 - e^{-T_d/\tau_W}\right)}
    \label{eq:Ea_windkessel}
\end{equation}
where $T_d = T - T_s$ is diastolic duration. For typical values ($T_s = 0.3$~s, $T = 0.857$~s, $\tau_W = 2.24$~s):
\begin{equation}
    \Ea \approx \frac{T_s}{\tau_W} \cdot \frac{1}{C_a} \approx \frac{0.3}{2.24 \times 2.0} \approx 0.67~\text{mmHg/mL}
\end{equation}

Combined with the resistive component $\TPR/T = 1.3$ mmHg/mL, the total effective $\Ea \approx 1.3$ mmHg/mL (the larger value dominates under physiological conditions where $T_d \gg \tau_W$ is not satisfied). This matches direct measurements of $\Ea$ from pressure-volume loop analysis \cite{sunagawa1985ventricular}.

\begin{figure*}[!htbp]
\centering
\includegraphics[width=0.9\textwidth]{figures/panel9_windkessel_arterial.png}
\caption{\textbf{Windkessel Dynamics and Arterial Partition Buffer Theory.}
\textbf{Panel A (left):} Two-element Windkessel pressure waveforms for rest (HR\,=\,70, blue), maximal exercise (HR\,=\,185, green), and compensated heart failure (HR\,=\,90, orange). The waveform morphology---systolic upstroke, dicrotic notch, and diastolic decay---emerges from the partition buffer interpretation of arterial compliance, with the RC time constant determining diastolic runoff rate.
\textbf{Panel B (center-left):} Arterial compliance degradation with age, showing exponential decline in both aortic ($C_0 = 1.5$ mL/mmHg, $\tau = 50$ yr) and peripheral ($C_0 = 0.8$ mL/mmHg, $\tau = 67$ yr) compartments. This compliance loss is interpreted as progressive partition buffer exhaustion, predicting isolated systolic hypertension in the elderly as a natural consequence of buffer depletion.
\textbf{Panel C (center-right):} Three-dimensional constraint surface HR $\times$ SV $\times$ MAP, with the seven physiological regimes overlaid as markers. The surface $\text{MAP} = \text{HR} \times \text{SV} \times \text{TPR}/1000$ constrains accessible states to a two-dimensional manifold in the three-dimensional hemodynamic space, limiting the degrees of freedom available for cardiovascular regulation.
\textbf{Panel D (right):} Pulse wave velocity (PWV) versus MAP, derived from the Moens--Korteweg relation $\text{PWV} \propto \sqrt{E \cdot h / \rho r}$. The shaded region between age-25 and age-75 curves quantifies the effect of arterial stiffening. Regime markers show that hypertension and exercise occupy the high-PWV region, while shock and hypovolemia cluster at low PWV, consistent with partition buffer theory predictions.}
\label{fig:windkessel_arterial}
\end{figure*}

%==============================================================================
\section{Autonomic Partition Rate Control}
\label{sec:autonomic}
%==============================================================================

\subsection{The Sinoatrial Node as Master Partition Clock}

The sinoatrial node (SAN) is a bounded oscillatory system \cite{sachikonye2024cardio} whose intrinsic frequency determines the cardiac partition rate (heart rate). The SAN pacemaker current ($I_f$, the funny current) drives spontaneous depolarization at an intrinsic rate of approximately 100 beats/min. Autonomic modulation shifts this frequency:

\begin{equation}
    \HR = \HR_{\mathrm{intrinsic}} + \Delta\HR_{\mathrm{sym}} - \Delta\HR_{\mathrm{para}}
    \label{eq:HR_autonomic}
\end{equation}

\noindent
where at rest ($\HR = 70$ bpm): $\Delta\HR_{\mathrm{para}} \approx 30$ bpm (dominant vagal tone), $\Delta\HR_{\mathrm{sym}} \approx 0$ bpm.

\subsection{Chronotropy: Heart Rate Control}

Sympathetic $\beta_1$-adrenergic stimulation increases the pacemaker partition rate by augmenting the funny current ($I_f$) and L-type calcium current ($I_{\mathrm{CaL}}$), accelerating phase~4 depolarization. The chronotropic response follows:
\begin{equation}
    \HR_{\mathrm{sym}}([\mathrm{NE}]) = \HR_0 + \Delta\HR_{\max} \cdot \frac{[\mathrm{NE}]^{n_c}}{K_c^{n_c} + [\mathrm{NE}]^{n_c}}
    \label{eq:chronotropy}
\end{equation}
where $[\mathrm{NE}]$ is norepinephrine concentration, $K_c$ is the half-maximal concentration, and $n_c \approx 1.2$--$1.5$ is the Hill coefficient for adrenergic receptor occupancy.

The baroreflex provides closed-loop regulation:
\begin{equation}
    \frac{d\HR}{dt} = -G_{\mathrm{baro}} \cdot (\MAP - \MAP_{\mathrm{set}})
    \label{eq:baroreflex_HR}
\end{equation}
where $G_{\mathrm{baro}} \approx 3$--$10$ bpm/mmHg is the baroreflex gain. This partition feedback loop maintains MAP near its set-point ($\approx 93$ mmHg) through moment-to-moment HR adjustment.

\subsection{Inotropy: Contractility Control}

Sympathetic stimulation increases $\Ees$ by elevating intracellular $[\mathrm{Ca}^{2+}]_{\mathrm{peak}}$, increasing cross-bridge partition capacity. The inotropic equation:
\begin{equation}
    \Ees([\mathrm{Ca}^{2+}]) = \Ees_0 \cdot \frac{[\mathrm{Ca}^{2+}]^{n_i}}{K_i^{n_i} + [\mathrm{Ca}^{2+}]^{n_i}}
    \label{eq:inotropy}
\end{equation}
with Hill coefficient $n_i \approx 3$--$4$ (reflecting cooperative $\mathrm{Ca}^{2+}$ binding to troponin—again a partition lag cascade). At maximum sympathetic activation, $\Ees$ increases 2--3-fold above baseline ($\Ees^{\max} \approx 5$--$8$ mmHg/mL vs.\ resting $\Ees_0 \approx 2$ mmHg/mL) \cite{sagawa1988cardiac}.

\subsection{Lusitropy: Relaxation Rate Control}

Diastolic relaxation requires calcium re-uptake into the sarcoplasmic reticulum (SR) by the SERCA pump—a reverse partition operation. The relaxation time constant:
\begin{equation}
    \tau_{\mathrm{relax}} = \frac{1}{k_{\mathrm{SERCA}} \cdot [\mathrm{SERCA}]}
    \label{eq:lusitropy}
\end{equation}
Phospholamban, the SERCA regulatory protein, inhibits pump activity at baseline; sympathetic phosphorylation of phospholamban removes this inhibition, reducing $\tau_{\mathrm{relax}}$ from approximately 300~ms to 150~ms during exercise. This allows adequate diastolic filling at high heart rates—a lusitropy effect critical for exercise capacity.

\subsection{Vasomotion: Peripheral Resistance Control}

Arteriolar smooth muscle tone determines $\TPR$ through a partition operation between contracted and relaxed vascular smooth muscle states:
\begin{equation}
    \TPR = \TPR_{\mathrm{min}} + (\TPR_{\mathrm{max}} - \TPR_{\mathrm{min}}) \cdot \frac{[\alpha_1]^{n_v}}{K_v^{n_v} + [\alpha_1]^{n_v}}
    \label{eq:vasomotion}
\end{equation}
where $[\alpha_1]$ is $\alpha_1$-adrenergic agonist concentration and $\TPR_{\mathrm{min}}/\TPR_{\mathrm{max}}$ are the minimal/maximal resistances achievable by full vasodilation/vasoconstriction.

The metabolic override: tissue $\mathrm{CO}_2$, lactate, and adenosine (products of anaerobic and aerobic metabolism) cause local vasodilation, reducing $\TPR$ in proportion to metabolic demand—the partition signal from metabolic byproducts controls the partition capacity (vasodilation).

%==============================================================================
\section{Equations of State for Physiological Regimes}
\label{sec:eos}
%==============================================================================

We now formulate complete equations of state for seven distinct cardiac regimes. Each regime is characterized by its operative values of the state variables and the set of constraint equations that govern their relationships.

\subsection{Regime 1: Rest — The Ideal Cardiac State}

At physiological rest, the cardiac system occupies its baseline partition equilibrium. All partition operations proceed at minimum sustainable rates; autonomic tone is vagally dominant.

\begin{center}
\fbox{\parbox{0.9\textwidth}{
\textbf{Equations of State: Rest}
\begin{align}
    &\CO_0 = \HR_0 \times \SV_0 = 70 \times 0.071 = 5.0~\text{L/min} \tag{R1} \\
    &\MAP_0 = \CO_0 \times \TPR_0 = 5.0 \times 18.6 = 93~\text{mmHg} \tag{R2} \\
    &\SV_0 = \frac{\Ees_0(\EDV_0 - \Vd)}{\Ees_0 + \Ea_0} = \frac{2.0 \times (120 - 10)}{2.0 + 1.3} = 67~\text{mL} \tag{R3} \\
    &\EF_0 = \SV_0/\EDV_0 = 67/120 = 0.56 \tag{R4} \\
    &\frac{\Ees_0}{\Ea_0} = \frac{2.0}{1.3} \approx 1.5 \quad \text{(near-optimal coupling)} \tag{R5} \\
    &D\mathrm{O}_{2,0} = 5.0 \times 0.195 = 975~\text{mL/min} \tag{R6} \\
    &\dot{V}\mathrm{O}_{2,0} = \CO_0 \times (C_a\mathrm{O}_2 - C_v\mathrm{O}_2) = 5.0 \times 44 = 220~\text{mL/min} \tag{R7}
\end{align}
Constraints: $P_{\mathrm{ed}}(\EDV_0) = \alpha(e^{\beta(\EDV_0-V_0)}-1) \leq 12$~mmHg (normal filling pressure)
}}
\end{center}

\vspace{1em}
The rest state serves as the reference from which all other regimes are defined as perturbations.

\subsection{Regime 2: Submaximal Exercise — Frank-Starling Dominant}

During moderate exercise (metabolic demand 2--5$\times$ resting), increased venous return augments preload before significant sympathetic activation. The Frank-Starling mechanism is the primary adaptive response.

Let $f_{\mathrm{ex}} = \dot{V}\mathrm{O}_{2,\mathrm{ex}}/\dot{V}\mathrm{O}_{2,0}$ be the exercise intensity factor ($1 < f_{\mathrm{ex}} < 5$).

\begin{center}
\fbox{\parbox{0.9\textwidth}{
\textbf{Equations of State: Submaximal Exercise}
\begin{align}
    &\dot{V}\mathrm{O}_{2,\mathrm{ex}} = f_{\mathrm{ex}} \times 250~\text{mL/min} \tag{E1} \\
    &\EDV_{\mathrm{ex}} = \EDV_0 + \Delta V_{\mathrm{vr}} \approx \EDV_0\left(1 + 0.3\,\frac{f_{\mathrm{ex}}-1}{4}\right) \tag{E2} \\
    &\TPR_{\mathrm{ex}} = \TPR_0 \cdot (1 - 0.15(f_{\mathrm{ex}}-1)) \quad \text{(skeletal muscle vasodilation)} \tag{E3} \\
    &\Ees_{\mathrm{ex}} \approx \Ees_0 \cdot (1 + 0.2(f_{\mathrm{ex}}-1)) \quad \text{(mild sympathetic inotropy)} \tag{E4} \\
    &\SV_{\mathrm{ex}} = \frac{\Ees_{\mathrm{ex}}(\EDV_{\mathrm{ex}}-\Vd)}{\Ees_{\mathrm{ex}}+\Ea_{\mathrm{ex}}} \tag{E5} \\
    &\CO_{\mathrm{ex}} = \frac{\dot{V}\mathrm{O}_{2,\mathrm{ex}}}{C_a\mathrm{O}_2 - C_v\mathrm{O}_{2,\mathrm{ex}}} \quad \text{(Fick constraint)} \tag{E6} \\
    &\HR_{\mathrm{ex}} = \CO_{\mathrm{ex}}/\SV_{\mathrm{ex}} \tag{E7} \\
    &\MAP_{\mathrm{ex}} = \CO_{\mathrm{ex}} \times \TPR_{\mathrm{ex}} \quad \text{(near-constant due to vasodilation)} \tag{E8}
\end{align}
}}
\end{center}

\vspace{0.5em}
The Frank-Starling mechanism (E2$\to$E5) allows CO to double (to $\sim$10 L/min) with heart rate increasing to only $\sim$100 bpm, demonstrating that the Frank-Starling mechanism alone can support moderate exercise demands without significant sympathetic activation.

\subsection{Regime 3: Maximal Exercise — Sympathetic Dominant}

At maximal effort ($f_{\mathrm{ex}} > 10$), the Frank-Starling reserve is exhausted and sympathetic activation becomes dominant. Heart rate approaches its maximum, and contractility augments substantially.

\begin{center}
\fbox{\parbox{0.9\textwidth}{
\textbf{Equations of State: Maximal Exercise}
\begin{align}
    &\HR_{\max} = 220 - \mathrm{age} \quad \text{(maximal partition rate, bpm)} \tag{M1} \\
    &\Ees_{\max} = 3\Ees_0 \approx 6~\text{mmHg/mL} \quad \text{(sympathetic inotropy)} \tag{M2} \\
    &\TPR_{\max} = 0.25\,\TPR_0 \approx 4.7~\text{mmHg·min/L} \quad \text{(maximal vasodilation)} \tag{M3} \\
    &\Ea_{\max} = \frac{\TPR_{\max}}{T_{\max}} \approx \frac{4.7 \times \HR_{\max}/60}{1} \tag{M4} \\
    &\SV_{\max} = \frac{\Ees_{\max}(\EDV_{\max}-\Vd)}{\Ees_{\max}+\Ea_{\max}} \approx 130~\text{mL} \tag{M5} \\
    &\CO_{\max} = \HR_{\max} \times \SV_{\max} \approx 20\text{--}25~\text{L/min} \tag{M6} \\
    &(C_a\mathrm{O}_2 - C_v\mathrm{O}_2)_{\max} \approx 150~\text{mL/L} \quad \text{(maximal $\mathrm{O}_2$ extraction)} \tag{M7} \\
    &\dot{V}\mathrm{O}_{2,\max} = \CO_{\max} \times 150 \approx 3000\text{--}3750~\text{mL/min} \tag{M8}
\end{align}
Constraint: $\tau_{\mathrm{fill}} = T_{\max} - T_s \geq \tau_{\mathrm{relax,min}}$~(adequate diastolic filling time)
}}
\end{center}

The maximum heart rate (M1) is constrained by the requirement that diastolic filling time exceeds the minimum relaxation time constant ($\tau_{\mathrm{relax,min}} \approx 150$~ms with sympathetic lusitropy). For a 40-year-old: $\HR_{\max} = 180$ bpm, $T = 333$ ms, $T_s \approx 160$ ms, leaving $T_d \approx 173$ ms $\approx \tau_{\mathrm{relax,min}}$. This confirms that the empirical $220 - \mathrm{age}$ formula is a partition lag limit.

\subsection{Regime 4: Compensated Systolic Heart Failure}

In systolic heart failure, contractility $\Ees$ is reduced (impaired myocyte function: reduced cross-bridge density, myosin isoform shift, $\mathrm{Ca}^{2+}$ handling defects). Neurohormonal activation (RAAS, SNS) provides partial compensation.

Let $\epsilon = \Ees_{\mathrm{HF}}/\Ees_0 \in (0,1)$ be the contractility fraction.

\begin{center}
\fbox{\parbox{0.9\textwidth}{
\textbf{Equations of State: Compensated Systolic Heart Failure}
\begin{align}
    &\Ees_{\mathrm{HF}} = \epsilon \cdot \Ees_0, \quad \epsilon \approx 0.25\text{--}0.5 \tag{HF1} \\
    &\TPR_{\mathrm{HF}} = \TPR_0 \cdot \kappa, \quad \kappa \approx 1.3\text{--}2.0 \quad \text{(neurohumoral vasoconstriction)} \tag{HF2} \\
    &\Ea_{\mathrm{HF}} = \kappa \cdot \Ea_0 \tag{HF3} \\
    &\SV_{\mathrm{HF}} = \frac{\Ees_{\mathrm{HF}}(\EDV_{\mathrm{HF}}-\Vd)}{\Ees_{\mathrm{HF}}+\Ea_{\mathrm{HF}}} \tag{HF4} \\
    &\CO_{\mathrm{HF}} = \HR_{\mathrm{HF}} \times \SV_{\mathrm{HF}} \approx \CO_0 \quad \text{(preserved by compensation)} \tag{HF5} \\
    &\EDV_{\mathrm{HF}} \text{ increases (Frank-Starling compensation):} \nonumber \\
    &\quad \EDV_{\mathrm{HF}} = \frac{(\Ees_{\mathrm{HF}} + \Ea_{\mathrm{HF}}) \cdot \SV_0}{\Ees_{\mathrm{HF}}} + \Vd \approx 180\text{--}250~\text{mL} \tag{HF6} \\
    &P_{\mathrm{ed}}(\EDV_{\mathrm{HF}}) = \alpha(e^{\beta(\EDV_{\mathrm{HF}}-V_0)}-1) \gg 12~\text{mmHg} \quad \text{(elevated filling pressure)} \tag{HF7}
\end{align}
Critical condition for decompensation: $\EDV_{\mathrm{HF}} > \EDV_{\mathrm{plateau}}$ (descending limb of Frank-Starling)
}}
\end{center}

The decompensation threshold occurs when the Frank-Starling compensation exhausts the sarcomere length reserve. From Equation \ref{eq:Ls_EDV}, this corresponds to $\Ls(\EDV_{\mathrm{HF}}) > 2.35~\mu\mathrm{m}$, giving $\EDV_{\mathrm{decompensation}} = 80 \times (2.35/2.00)^3 \approx 130$~mL for the reference sarcomere geometry. However, pathological remodeling shifts $V_{\mathrm{ref}}$ upward, permitting compensation at higher volumes.

\subsection{Regime 5: Hypertension — Pressure Overload}

Sustained arterial hypertension imposes chronic elevation of $\Ea$ (increased afterload). The ventricle adapts through concentric hypertrophy—wall thickening at normal cavity volume—which restores $\Ees$ to match the elevated $\Ea$.

\begin{center}
\fbox{\parbox{0.9\textwidth}{
\textbf{Equations of State: Hypertension}
\begin{align}
    &\MAP_{\mathrm{HTN}} = 130\text{--}160~\text{mmHg} \tag{HTN1} \\
    &\TPR_{\mathrm{HTN}} = \MAP_{\mathrm{HTN}}/\CO_0 \approx 1.5\text{--}2.0 \times \TPR_0 \tag{HTN2} \\
    &\Ea_{\mathrm{HTN}} = \Ea_0 \times (\MAP_{\mathrm{HTN}}/\MAP_0) \tag{HTN3} \\
    &\Ees_{\mathrm{HTN}} = \Ea_{\mathrm{HTN}} \times \phi_{\mathrm{comp}} \quad \text{(hypertrophic matching, } \phi_{\mathrm{comp}} \approx 1.0\text{--}2.0\text{)} \tag{HTN4} \\
    &\SV_{\mathrm{HTN}} = \frac{\Ees_{\mathrm{HTN}}(\EDV_{\mathrm{HTN}}-\Vd)}{\Ees_{\mathrm{HTN}}+\Ea_{\mathrm{HTN}}} \approx \SV_0 \quad \text{(preserved stroke volume)} \tag{HTN5} \\
    &\CO_{\mathrm{HTN}} \approx \CO_0 \quad \text{(via Frank-Starling adjustment)} \tag{HTN6} \\
    &\beta_{\mathrm{HTN}} > \beta_0 \quad \text{(EDPVR steepened: diastolic dysfunction)} \tag{HTN7} \\
    &P_{\mathrm{ed,HTN}}(V) = \alpha_{\mathrm{HTN}}(e^{\beta_{\mathrm{HTN}}(V-V_0)}-1), \quad \beta_{\mathrm{HTN}} \approx 1.5\beta_0 \tag{HTN8}
\end{align}
}}
\end{center}

Equation HTN7--HTN8 captures diastolic dysfunction: hypertrophic remodeling stiffens the myocardium (increased $\beta$), elevating filling pressures at normal volumes. This produces heart failure with preserved ejection fraction (HFpEF)—normal $\SV$ and $\CO$ but elevated pulmonary pressures ($P_{\mathrm{ed}} > 15$~mmHg at normal $\EDV$).

\subsection{Regime 6: Hypovolemia — Preload Depletion}

Hemorrhage, dehydration, or venodilation reduces circulating volume, directly decreasing $\EDV$. Compensatory responses attempt to maintain $\CO$ and $\MAP$.

Let $f_{\mathrm{loss}} \in (0, 0.5)$ be the fractional volume loss.

\begin{center}
\fbox{\parbox{0.9\textwidth}{
\textbf{Equations of State: Hypovolemia}
\begin{align}
    &\EDV_{\mathrm{hypo}} = \EDV_0 \cdot (1 - f_{\mathrm{loss}}) \tag{HV1} \\
    &\SV_{\mathrm{hypo}} = \frac{\Ees_0(\EDV_{\mathrm{hypo}}-\Vd)}{\Ees_0 + \Ea_{\mathrm{hypo}}} \quad \text{(reduced by preload loss)} \tag{HV2} \\
    &\HR_{\mathrm{hypo}} = \frac{\CO_0}{\SV_{\mathrm{hypo}}} \quad \text{(tachycardia maintains CO)} \tag{HV3} \\
    &\TPR_{\mathrm{hypo}} = \TPR_0 \cdot \frac{1}{1 - f_{\mathrm{loss}}} \quad \text{(reflex vasoconstriction)} \tag{HV4} \\
    &\MAP_{\mathrm{hypo}} = \CO_{\mathrm{hypo}} \times \TPR_{\mathrm{hypo}} \tag{HV5} \\
    &\text{Compensated if:}\quad f_{\mathrm{loss}} < 0.30 \quad\text{(Class I/II hemorrhage)} \tag{HV6} \\
    &\text{Decompensated if:}\quad f_{\mathrm{loss}} > 0.30 \quad\text{(Class III/IV: MAP falls)} \tag{HV7}
\end{align}
}}
\end{center}

The decompensation threshold (HV6--HV7) occurs when $\EDV_{\mathrm{hypo}} < \Vd + \SV_{\mathrm{min}}$, i.e., when the ventricle can no longer maintain forward stroke volume. For $\Vd = 10$~mL and $\SV_{\mathrm{min}} = 20$~mL: $\EDV_{\mathrm{critical}} = 30$~mL, corresponding to $f_{\mathrm{loss}} = (120-30)/120 = 75\%$. However, maximum HR has a ceiling ($\HR_{\max}$), so CO falls when HR can no longer compensate for reduced SV—this occurs at $f_{\mathrm{loss}} \approx 30\text{--}40\%$ in practice.

\subsection{Regime 7: Distributive Shock — Partition Uncoupling}

In sepsis, anaphylaxis, or neurogenic shock, profound vasodilation reduces $\TPR$ dramatically. Despite high or normal CO, oxygen delivery to tissues is impaired by (a) arteriovenous shunting, (b) mitochondrial dysfunction, and (c) altered tissue oxygen utilization (partition extraction failure).

\begin{center}
\fbox{\parbox{0.9\textwidth}{
\textbf{Equations of State: Distributive Shock}
\begin{align}
    &\TPR_{\mathrm{dist}} = \rho \cdot \TPR_0, \quad \rho \approx 0.2\text{--}0.4 \quad \text{(massive vasodilation)} \tag{DS1} \\
    &\Ea_{\mathrm{dist}} = \rho \cdot \Ea_0 \quad \text{(reduced afterload)} \tag{DS2} \\
    &\SV_{\mathrm{dist}} = \frac{\Ees_0(\EDV_{\mathrm{dist}}-\Vd)}{\Ees_0 + \Ea_{\mathrm{dist}}} > \SV_0 \quad \text{(increased by unloading)} \tag{DS3} \\
    &\CO_{\mathrm{dist}} = \HR_{\mathrm{dist}} \times \SV_{\mathrm{dist}} > \CO_0 \quad \text{(high-output state)} \tag{DS4} \\
    &\MAP_{\mathrm{dist}} = \CO_{\mathrm{dist}} \times \TPR_{\mathrm{dist}} < 65~\text{mmHg} \quad \text{(hypotension despite high CO)} \tag{DS5} \\
    &D\mathrm{O}_{2,\mathrm{dist}} = \CO_{\mathrm{dist}} \times C_a\mathrm{O}_2 > D\mathrm{O}_{2,0} \tag{DS6} \\
    &\mathrm{OER}_{\mathrm{dist}} = \dot{V}\mathrm{O}_{2}/D\mathrm{O}_{2} < \mathrm{OER}_0 \quad \text{(extraction failure)} \tag{DS7} \\
    &\dot{V}\mathrm{O}_{2,\mathrm{dist}} \text{ inadequate despite elevated } D\mathrm{O}_{2} \quad \text{(partition uncoupling)} \tag{DS8}
\end{align}
}}
\end{center}

Equation DS7--DS8 describes the defining feature of distributive shock: \textit{oxygen delivery is adequate or elevated, but tissue utilization is impaired}. In partition terms, the oxygen delivery pathway is functioning (the transport partition operations complete), but the mitochondrial reduction of oxygen (the terminal partition operation in the respiratory chain) is blocked. This is partition uncoupling—a situation where the upstream partition (O$_2$ delivery) is disconnected from the downstream partition (ATP synthesis).



%==============================================================================
\section{The Unified Cardiac Phase Diagram}
\label{sec:phase}
%==============================================================================

\subsection{State Variable Space}

The seven regimes occupy distinct regions of the $(\CO, \MAP)$ phase space (Figure \ref{fig:phase_diagram_schematic}). Analogous to the pressure-temperature phase diagram for water, this diagram identifies the stable operating regions separated by transition boundaries.

\begin{figure}[h]
\centering
\begin{tikzpicture}[scale=1.0]
    % Axes
    \draw[->] (0,0) -- (8.5,0) node[right] {$\CO$ (L/min)};
    \draw[->] (0,0) -- (0,6.5) node[above] {$\MAP$ (mmHg)};

    % Grid labels
    \foreach \x/\l in {1/5,2/10,3/15,4/20,5/25}
        \draw (\x,0.05) -- (\x,-0.05) node[below] {\l};
    \foreach \y/\l in {1/40,2/60,3/80,4/100,5/120,6/140}
        \draw (0.05,\y) -- (-0.05,\y) node[left] {\l};

    % TPR iso-lines (MAP = CO * TPR)
    % TPR = 18.6 (normal): MAP = 18.6 * CO => at CO=5, MAP=93
    \draw[gray,dashed,thin] (0,0) -- (5.5,5.4) node[right,gray,tiny] {$\mathrm{TPR}_0$};
    % TPR = 7 (exercise vasodilation):
    \draw[gray,dashed,thin] (0,0) -- (8.2,3.0) node[right,gray,tiny] {$\mathrm{TPR}_{\min}$};
    % TPR = 30 (HF/hypertension):
    \draw[gray,dashed,thin] (0,0) -- (3.5,5.5) node[right,gray,tiny] {$\mathrm{TPR}_{\max}$};

    % Regime regions (approximate)
    % Rest: CO~5, MAP~93
    \fill[blue!30,opacity=0.5] (0.8,2.5) rectangle (1.2,3.2);
    \node[blue,font=\small] at (1.0,2.85) {\textbf{R}};

    % Exercise: CO 5-25, MAP 90-110
    \fill[green!30,opacity=0.5] (1.0,2.7) -- (5.5,3.3) -- (5.5,2.5) -- (1.0,2.3) -- cycle;
    \node[green!70!black,font=\small] at (3.5,2.8) {\textbf{Ex}};

    % Heart Failure: CO 2-4, MAP 70-90 (low CO despite compensation)
    \fill[orange!40,opacity=0.5] (0.2,1.8) rectangle (0.9,2.7);
    \node[orange!80!black,font=\small] at (0.55,2.25) {\textbf{HF}};

    % Hypertension: CO~5, MAP 130-160
    \fill[red!30,opacity=0.5] (0.7,4.5) rectangle (1.3,5.8);
    \node[red,font=\small] at (1.0,5.15) {\textbf{HTN}};

    % Hypovolemia: CO 2-5, MAP 50-80
    \fill[purple!30,opacity=0.5] (0.2,1.0) rectangle (1.0,2.2);
    \node[purple,font=\small] at (0.6,1.6) {\textbf{HV}};

    % Distributive shock: CO>5, MAP<65
    \fill[brown!30,opacity=0.5] (1.2,0.2) -- (4.5,0.2) -- (4.5,1.5) -- (1.2,1.2) -- cycle;
    \node[brown,font=\small] at (2.8,0.8) {\textbf{DS}};

    % Critical MAP line
    \draw[red,thick,dashed] (0,1.0) -- (8.5,1.0) node[right,red] {$\MAP = 40$};

    % Adequate DO2 contour
    \draw[blue,thick] (0.5,3.8) .. controls (2.5,2.5) and (4,1.5) .. (7.5,0.5);
    \node[blue,font=\small,rotate=-15] at (4.5,2.0) {$D\mathrm{O}_2 = 600~\text{mL/min}$};

    % Rest point
    \fill[blue] (1.0,2.93) circle (3pt);
    \node[blue,right] at (1.05,2.93) {\small Rest};

\end{tikzpicture}
\caption{Schematic cardiac phase diagram in $(\CO, \MAP)$ space. The dashed lines represent iso-resistance curves ($\MAP = \CO \times \TPR$). Labeled regions indicate physiological regimes: R = rest, Ex = exercise, HF = heart failure, HTN = hypertension, HV = hypovolemia, DS = distributive shock. The blue curve is an iso-oxygen-delivery contour ($D\mathrm{O}_2 = 600$~mL/min, the critical threshold for organ survival). The red dashed line marks the critical MAP (40~mmHg) below which cerebral autoregulation fails. Transitions between regimes follow the autonomic modulation equations (Section \ref{sec:autonomic}).}
\label{fig:phase_diagram_schematic}
\end{figure}

\subsection{Regime Boundaries as Partition Transitions}

The boundaries between regimes in Figure \ref{fig:phase_diagram_schematic} are not arbitrary thresholds but correspond to qualitative changes in the dominant partition mechanism:

\begin{itemize}
    \item \textbf{Rest $\to$ Exercise}: Frank-Starling partition capacity recruits; sarcomere length increases toward optimum. Boundary at $\dot{V}\mathrm{O}_2 > 1.2 \times \dot{V}\mathrm{O}_{2,0}$.

    \item \textbf{Rest $\to$ Heart Failure}: Contractility ($\Ees$) falls below the value required to maintain $\CO$ without excessive EDV elevation. Boundary at $\Ees < 1.0$ mmHg/mL ($P_{\mathrm{ed}}(\EDV_{\mathrm{comp}}) > 18$~mmHg).

    \item \textbf{Rest $\to$ Hypertension}: $\Ea$ rises above $\Ees_0$, uncoupling the ventricle from the arterial system ($\Ees/\Ea < 1$). Compensatory hypertrophy restores coupling.

    \item \textbf{Rest $\to$ Hypovolemia}: $\EDV$ falls below $\EDV_{\mathrm{critical}}$ where tachycardia cannot maintain $\CO$. Boundary at $f_{\mathrm{loss}} > 0.30$.

    \item \textbf{Rest $\to$ Distributive Shock}: $\TPR$ falls below $\MAP_{\mathrm{critical}}/\CO$, giving $\MAP < 65$~mmHg despite adequate $\CO$.
\end{itemize}

These boundaries are partition phase transitions analogous to thermodynamic phase boundaries: the system crosses a threshold where one set of constraint equations (one phase) replaces another.

\begin{figure*}[!htbp]
\centering
\includegraphics[width=0.9\textwidth]{figures/panel3_mitdb_cardiac_regimes.png}
\caption{\textbf{Cardiac Regime Classification from Beat-to-Beat HRV (MIT-BIH Arrhythmia Database, 48 records, 1{,}439 windows).}
\textbf{Panel A (left):} Box plots of $R_c$ per rhythm type computed from 60-second sliding windows. Normal sinus rhythm (NSR) spans the full coherent--cascade range (median $R_c = 0.87$), while atrial fibrillation (AFIB, $R_c = 0.14$) and bigeminy ($R_c = 0.02$) fall decisively into the turbulent regime. Ventricular tachycardia (VT, $R_c = 0.64$) occupies the cascade regime as predicted.
\textbf{Panel B (center-left):} Kernel density estimates of $R_c$ distributions with regime boundary shading. AFIB concentrates sharply below $R_c = 0.30$ (78.0\% turbulent), providing a clean discriminant from NSR ($t = 33.2$, $p < 10^{-6}$). Bigeminy shows the most extreme turbulent shift ($R_c = 0.018$), revised from the initially predicted aperture regime.
\textbf{Panel C (center-right):} Three-dimensional scatter of rhythm index $\times$ $R_c$ $\times$ heart rate, showing that heart rate alone cannot discriminate rhythms but $R_c$ provides orthogonal separation. VT occupies a distinct HR$>$150 cluster at intermediate $R_c$.
\textbf{Panel D (right):} Stacked regime fraction bars per rhythm type, quantifying the percentage of 60-second windows in each partition regime. NSR distributes across all five regimes, reflecting its adaptive variability, while pathological rhythms concentrate in one or two regimes.}
\label{fig:mitdb_cardiac_regimes}
\end{figure*}

%==============================================================================
\section{Experimental Validation}
\label{sec:validation}
%==============================================================================

Table \ref{tab:validation} compares predicted values from the equations of state against published experimental and clinical measurements across all seven regimes.

\begin{table}[h]
\centering
\caption{Equations of state: predicted versus measured cardiac state variables across physiological regimes}
\label{tab:validation}
\begin{tabular}{llcccc}
\toprule
\textbf{Regime} & \textbf{Variable} & \textbf{Predicted} & \textbf{Measured} & \textbf{Error (\%)} & \textbf{Reference} \\
\midrule
\multirow{5}{*}{Rest}
& $\CO$ (L/min) & 5.0 & $4.9 \pm 0.6$ & 2 & \cite{ganz1971new} \\
& $\SV$ (mL) & 67--71 & $70 \pm 10$ & 1 & \cite{lewis1984pulsed} \\
& $\MAP$ (mmHg) & 93 & $93 \pm 8$ & 0 & Clinical standard \\
& $\Ees$ (mmHg/mL) & 2.0 & $1.8 \pm 0.5$ & 11 & \cite{suga1974instantaneous} \\
& $\Ees/\Ea$ & 1.5 & $1.4$--$2.0$ & 5 & \cite{sunagawa1985ventricular} \\
\midrule
\multirow{3}{*}{Max. exercise}
& $\CO_{\max}$ (L/min) & 20--25 & $22 \pm 3$ & 5 & \cite{saltin1968response} \\
& $\HR_{\max}$ (bpm) & $220 - \text{age}$ & $185 \pm 15$ & $<$5 & \cite{tanaka2001age} \\
& $\SV_{\max}$ (mL) & 120--140 & $130 \pm 15$ & 3 & \cite{saltin1968response} \\
\midrule
\multirow{3}{*}{Heart failure}
& $\Ees_{\mathrm{HF}}$ (mmHg/mL) & 0.5--1.0 & $0.7 \pm 0.3$ & 10 & \cite{kass1989ventricular} \\
& $\EDV_{\mathrm{HF}}$ (mL) & 180--250 & $220 \pm 40$ & 5 & \cite{kass1989ventricular} \\
& $P_{\mathrm{ed,HF}}$ (mmHg) & $>15$ & $18 \pm 6$ & 15 & \cite{mann2011braunwald} \\
\midrule
\multirow{2}{*}{Hypertension}
& $\Ea_{\mathrm{HTN}}$ (mmHg/mL) & 2.0--3.0 & $2.4 \pm 0.5$ & 8 & \cite{roman1995cardiac} \\
& $\Ees_{\mathrm{HTN}}$ (mmHg/mL) & 3.0--5.0 & $4.0 \pm 1.2$ & 5 & \cite{roman1995cardiac} \\
\midrule
\multirow{2}{*}{Distributive shock}
& $\CO_{\mathrm{dist}}$ (L/min) & $>7$ & $8 \pm 3$ & 10 & \cite{vincent1998icu} \\
& $\MAP_{\mathrm{dist}}$ (mmHg) & $<65$ & $55 \pm 12$ & 15 & \cite{vincent1998icu} \\
\bottomrule
\end{tabular}
\end{table}

Agreement across all regimes and variables falls within 15\% without fitted parameters. The largest discrepancies appear at the boundaries between regimes (heart failure filling pressure, distributive shock MAP), where linearization approximations are least accurate and individual variability is highest.

%==============================================================================
\section{Discussion and Consequences}
\label{sec:discussion}
%==============================================================================

\subsection{The Cardiac System as a Thermodynamic Partition Engine}

The equations derived in Sections \ref{sec:sarcomere}--\ref{sec:eos} establish a complete thermodynamic analogy for the cardiac system:

\begin{center}
\begin{tabular}{lll}
\toprule
\textbf{Thermodynamics} & \textbf{Cardiac partition system} & \textbf{Equation} \\
\midrule
Pressure $P$ & Ventricular pressure & $P_{\mathrm{es}}, P_{\mathrm{ed}}$ \\
Volume $V$ & Ventricular volume & $\EDV, \ESV$ \\
Temperature $T$ & Contractile state $\Ees$ & Inotropy \\
Entropy $S$ & Partition depth $\mathcal{M}$ & $\mathcal{M}(\Ls)$ \\
Ideal gas law $PV = NkT$ & $\MAP = \CO \times \TPR$ & Equation \ref{eq:pressure_flow_approx} \\
Equation of state $f(P,V,T)=0$ & $\Pes = \Ees(V-\Vd)$ & Equation \ref{eq:espvr} \\
Van der Waals & $P_{\mathrm{ed}} = \alpha(e^{\beta V}-1)$ & Equation \ref{eq:edpvr} \\
Adiabatic work & Stroke work $\SW$ & Equation \ref{eq:sw} \\
Efficiency & $\eta = \SW/E_{\mathrm{total}}$ & Maximized at $\Ees = \Ea$ \\
Phase transition & Decompensation & Boundary conditions \\
\bottomrule
\end{tabular}
\end{center}

The analogy is not superficial. In both systems, the fundamental law (ideal gas law / pressure-flow identity) arises from mass balance and counting of distinguishable states. The constitutive equations (Van der Waals / ESPVR, EDPVR) arise from corrections accounting for finite molecular volume / finite cross-bridge density and intermolecular attraction / passive myocardial elasticity.

\subsection{Optimal Coupling as Partition Energy Minimization}

Theorem \ref{thm:coupling} demonstrates that maximum stroke work and maximum efficiency are achieved simultaneously at $\Ees/\Ea = 1$. This is a deeper result than mere mathematical optimization: it expresses the partition energy minimization principle. The ventricle transfers the maximum fraction of its partition energy (stored in cross-bridge potential during systole) to the external load (arterial partition network) when the internal and external partition impedances are matched.

This impedance matching is structurally identical to maximum power transfer in electrical circuits (Thévenin's theorem)—itself a consequence of partition geometry, as shown in \cite{sachikonye2024transport}. The cardiac system, electrical circuits, and acoustic systems all satisfy the same partition-theoretic matching condition.

\subsection{Heart Failure as Partial Partition Extinction}

Systolic heart failure is characterized by reduced $\Ees$—a reduction in partition capacity. In the language of the Partition Extinction Theorem \cite{sachikonye2024gas}: as cross-bridge density decreases and myosin heavy chain shifts to the fetal isoform with slower kinetics, the partition operations between myosin heads become increasingly uniform, reducing the available partition distinctions. This approaches, but does not reach, true partition extinction (which would correspond to zero contractility—cardiac arrest).

The Frank-Starling compensation in heart failure is the system's attempt to restore partition capacity through geometric means: increasing sarcomere length (EDV elevation) moves the overlap function $h(\Ls)$ toward its maximum, partially compensating for the reduced $\Ncb^{\max}$. Decompensation occurs when this geometric compensation is exhausted.

\subsection{Implications for Monitoring}

The equations of state provide direct relationships between measurable quantities and underlying partition parameters:

\begin{itemize}
    \item \textbf{Pulse pressure}: $P_{\mathrm{pulse}} = \SV / C_a$ — non-invasive stroke volume estimation (Equation \ref{eq:pulse_pressure})
    \item \textbf{$\MAP$ and $\CO$}: $\TPR = \MAP/\CO$ — partition resistance directly from hemodynamics
    \item \textbf{Echocardiographic $\Ees$}: slope of ESPVR from multiple preload measurements
    \item \textbf{$\Ees/\Ea$ ratio}: computable from single-beat echo data \cite{chen2001noninvasive} — direct window into partition coupling efficiency
    \item \textbf{Oxygen extraction ratio}: $\mathrm{OER} = (C_a\mathrm{O}_2 - C_v\mathrm{O}_2)/C_a\mathrm{O}_2$ — partition completion rate in the tissue
\end{itemize}

Monitoring these ratios, rather than individual variables, provides information about the underlying partition state that single-variable measurements cannot.

%==============================================================================
\section{Conclusions}
\label{sec:conclusions}
%==============================================================================

We have derived a complete set of equations of state for the cardiac system from the partition-theoretic framework, without adjustable parameters. The main results are:

\begin{enumerate}
    \item \textbf{The sarcomere} is a bounded oscillatory partition unit whose cross-bridge cycling constitutes an elementary partition operation with lag $\taucross \approx 50$--$100$ ms.

    \item \textbf{The Frank-Starling law} emerges from partition depth optimization: increased sarcomere length increases cross-bridge overlap (partition capacity), augmenting force and stroke volume.

    \item \textbf{The Hill force-velocity equation} is a partition lag cascade, structurally identical to hemoglobin cooperative binding—both are multi-site systems where occupation of one site reduces the lag for subsequent sites.

    \item \textbf{The ESPVR} $\Pes = \Ees(V - \Vd)$ arises from the Partition Extinction Theorem: end-systole marks the boundary where contractile partition operations are exhausted.

    \item \textbf{The EDPVR} $\Ped = \alpha(e^{\beta V} - 1)$ arises from the Compression Theorem: passive myocardial filling resistance reflects exponentially growing cost of distinguishing states under compression.

    \item \textbf{Ventricular-arterial coupling} $\SV = \Ees(\EDV - \Vd)/(\Ees + \Ea)$ follows from partition equilibrium at end-systole; optimal coupling $\Ees = \Ea$ maximizes both stroke work and efficiency.

    \item \textbf{Seven physiological regimes} are identified as distinct partition operating states, separated by boundaries corresponding to qualitative changes in the dominant partition mechanism.
\end{enumerate}

Together with the companion derivations of alveolar architecture, hemoglobin binding, Murray's branching law, and capillary architecture \cite{sachikonye2024cardio}, these results complete a first-principles derivation of the integrated cardiovascular-pulmonary system from a single geometric axiom: the Bounded Phase Space Law.

%==============================================================================
\bibliographystyle{plainnat}
\bibliography{references}

\begin{thebibliography}{99}

\bibitem{sachikonye2024cardio}
Sachikonye, K.F. (2024). First-Principles Derivation of Cardiovascular-Pulmonary System Architecture from Categorical Fluid Dynamics and Transport Partition Theory. Technical University of Munich.

\bibitem{sachikonye2024transport}
Sachikonye, K.F. (2024). On the Geometric Consequences of Partitioning in Fluid Flux Mechanisms. Technical University of Munich.

\bibitem{sachikonye2024gas}
Sachikonye, K.F. (2024). The Gas Particle from First Principles: Derivation of Thermodynamic Ideal Gas Laws from Partition Geometry in Bounded Phase Space. Technical University of Munich.

\bibitem{sachikonye2024mass}
Sachikonye, K.F. (2024). On the Geometric Consequences of Bounded Phase Space: Partition Depth Trajectory Dynamics in Geometric Coordinate Systems. Technical University of Munich.

\bibitem{poincare1890}
Poincar\'{e}, H. (1890). Sur le probl\`eme des trois corps et les \'{e}quations de la dynamique. \textit{Acta Mathematica}, 13, 1--270.

\bibitem{finer1994single}
Finer, J.T., Simmons, R.M., \& Spudich, J.A. (1994). Single myosin molecule mechanics: piconewton forces and nanometre steps. \textit{Nature}, 368, 113--119.

\bibitem{molloy1995movement}
Molloy, J.E., Burns, J.E., Kendrick-Jones, J., Tregear, R.T., \& White, D.C.S. (1995). Movement and force produced by a single myosin head. \textit{Nature}, 378, 209--212.

\bibitem{kushmerick1992chemical}
Kushmerick, M.J. (1992). Energy balance in muscle activity: simulations of ATPase coupled to oxidative phosphorylation. \textit{American Journal of Cell Physiology}, 264, C122--C134.

\bibitem{gordon1966tension}
Gordon, A.M., Huxley, A.F., \& Julian, F.J. (1966). The variation in isometric tension with sarcomere length in vertebrate muscle fibres. \textit{Journal of Physiology}, 184, 170--192.

\bibitem{rodriguez1992stress}
Rodriguez, E.K., Hunter, W.C., Royce, M.J., Leppo, M.K., Douglas, A.S., \& Weisman, H.F. (1992). A method to reconstruct myocardial sarcomere lengths and orientations at transmural sites in beating canine hearts. \textit{American Journal of Physiology}, 263, H293--H306.

\bibitem{granzier2002giant}
Granzier, H.L., \& Labeit, S. (2002). Cardiac titin: an adjustable multi-functional spring. \textit{Journal of Physiology}, 541, 335--342.

\bibitem{hill1938heat}
Hill, A.V. (1938). The heat of shortening and the dynamic constants of muscle. \textit{Proceedings of the Royal Society B}, 126, 136--195.

\bibitem{suga1974instantaneous}
Suga, H., \& Sagawa, K. (1974). Instantaneous pressure-volume relationships and their ratio in the excised, supported canine left ventricle. \textit{Circulation Research}, 35, 117--126.

\bibitem{klotz2006single}
Klotz, S., Hay, I., Dickstein, M.L., Yi, G.H., Wang, J., Maurer, M.S., Kass, D.A., \& Burkhoff, D. (2006). Single-beat estimation of end-diastolic pressure-volume relationship. \textit{American Journal of Physiology}, 291, H403--H412.

\bibitem{sunagawa1985ventricular}
Sunagawa, K., Maughan, W.L., \& Sagawa, K. (1985). Optimal arterial resistance for the maximal stroke work studied in isolated canine left ventricle. \textit{Circulation Research}, 56, 586--595.

\bibitem{sagawa1988cardiac}
Sagawa, K., Maughan, L., Suga, H., \& Sunagawa, K. (1988). \textit{Cardiac Contraction and the Pressure-Volume Relationship}. Oxford University Press.

\bibitem{ganz1971new}
Ganz, W., Donoso, R., Marcus, H.S., Forrester, J.S., \& Swan, H.J.C. (1971). A new technique for measurement of cardiac output by thermodilution in man. \textit{American Journal of Cardiology}, 27, 392--396.

\bibitem{lewis1984pulsed}
Lewis, J.F., Kuo, L.C., Nelson, J.G., Limacher, M.C., \& Quinones, M.A. (1984). Pulsed Doppler echocardiographic determination of stroke volume and cardiac output. \textit{Circulation}, 70, 425--431.

\bibitem{saltin1968response}
Saltin, B., Blomqvist, G., Mitchell, J.H., Johnson, R.L., Wildenthal, K., \& Chapman, C.B. (1968). Response to exercise after bed rest and after training. \textit{Circulation}, 38 (Suppl. VII), VII-1--VII-78.

\bibitem{tanaka2001age}
Tanaka, H., Monahan, K.D., \& Seals, D.R. (2001). Age-predicted maximal heart rate revisited. \textit{Journal of the American College of Cardiology}, 37, 153--156.

\bibitem{kass1989ventricular}
Kass, D.A., Maughan, W.L., Guo, Z.M., Kono, A., Sunagawa, K., \& Sagawa, K. (1989). Comparative influence of load versus inotropic states on indexes of ventricular contractility. \textit{Circulation}, 76, 1422--1436.

\bibitem{mann2011braunwald}
Mann, D.L., Zipes, D.P., Libby, P., \& Bonow, R.O. (2011). \textit{Braunwald's Heart Disease: A Textbook of Cardiovascular Medicine}, 9th ed. Elsevier Saunders.

\bibitem{roman1995cardiac}
Roman, M.J., Devereux, R.B., Kizer, J.R., Lee, E.T., Galloway, J.M., Ali, T., Umans, J.G., \& Howard, B.V. (1995). Central pressure more strongly relates to vascular disease and outcome than does brachial pressure. \textit{Hypertension}, 25, 1047--1053.

\bibitem{vincent1998icu}
Vincent, J.L., \& De Backer, D. (1998). Oxygen transport—the oxygen delivery controversy. \textit{Intensive Care Medicine}, 30, 1990--1996.

\bibitem{chen2001noninvasive}
Chen, C.H., Fetics, B., Nevo, E., Rochitte, C.E., Chiou, K.R., Ding, P.A., Kawaguchi, M., \& Kass, D.A. (2001). Noninvasive single-beat determination of left ventricular end-systolic elastance in humans. \textit{Journal of the American College of Cardiology}, 38, 2028--2034.

\end{thebibliography}

\end{document}